\section{School-policy benchmarks and institutional sensitivity}\label{sec:school-sensitivity}
\paragraph{A predictor that simulates the allocation}
For 2022 and 2023 we used completed reports and priority tiers while withholding the within-tier order. Simulating DA to predict admission risk has precedent in \citet[online appendix, Sec.~F.1]{smartplatforms}. Here we hold all submitted lists fixed, draw an independent uniform order within each course and tier, and run DA on the full lists. Repeat 64 times. For each pupil keep the five discretionary programs they received most often, breaking frequency ties by submitted order and retaining the entitlement. We evaluate the resulting menus at the published order. This benchmark requires completed reports and uses no historical offers in selection.

Across eight independent 64-draw batches per year, this rule beat the rank-informed predictor on both recorded-offer coverage and placement in every comparison. It showed 494--541 more historical offers and placed 25--56 more pupils, but gave 106--163 fewer pupils their first choice. Compared instead with the rank-blind predictor, it still covered more historical offers and placed fewer people in every batch. Simulation precision matters: at 16 draws, only five of the 16 comparisons improved on rank-informed on both outcomes. In the first 64-draw batch there was still a zero-frequency tie at the fifth-slot boundary for 32,982 of 42,123 longer lists in 2022 and 35,800 of 43,409 in 2023. The fallback to submitted rank is part of the rule. The eight batches vary the estimated menus within each fixed cohort.

We then fixed the comparison, including all three simulation lengths and eight batches, before generating the new rule's 2024 outcomes. The cohort and the original predictors' outcomes were already known. At 64 draws the historical coverage advantage persisted in all eight batches ($+447$ to $+483$ offers), as did a much smaller placement advantage ($+2$ to $+24$ pupils). First choices fell by 105--123. At 16 draws \emph{none} of the eight batches placed more pupils; at 32 only three did. We did not extend the simulated-lottery test to 2025, the year the ministry replaced random tie-breaking.

Pooling all seats simplifies the ministry's allocation. A sensitivity that first honors continuity and gives designated pupils priority for earmarked vacancies changes the year-specific net placement losses to 469 and 516. Neither version models all reserves or dynamic family priorities; overlapping reserve categories in particular can change who is admitted \citep{doganreserves}. We also tested a deliberately generous seat count, $v_j+h_j$, to allow for the possibility that vacancies exclude continuing pupils. This adds 174,270--225,476 seats per year and reduces full-list agreement with historical outcomes to 82.01--85.68\%. With the original menus frozen the placement differences remain negative in all years. Recomputing demand-to-capacity features and choosing new menus under the same 2019 coefficients yields 337--498 fewer placements and 839--1,593 more historical offers shown per year. The two seat conventions are scenarios; neither recovers the ministry's full mechanism. Only the two frozen predictors' historical-offer and placement directions were specified before acquiring the validation cohorts.

\paragraph{List length}
Keeping both fitted scores fixed, we repeated the comparison at one, three, five, eight and ten discretionary slots in 2022 and 2023. In all ten year--quota combinations rank-informed covered more historical offers and placed fewer pupils in replay. The size varied sharply. At one slot, historical coverage rose by 37,861 and 42,580, while placement fell by 18,907 and 17,922. At ten slots, coverage rose by only 34 and 70 and placement fell by 48 and 51. Once all reported options fit, the policies give the same allocation by construction.

At the five-slot quota we also varied how ties are resolved within the original priority tiers in the 2022--23 replays. We drew new independent course-specific lotteries and, separately, a common applicant lottery across courses. We used 32 paired draws per ordering scheme and year, recalculating the full-list reference for each draw. The rank-informed policy raised full-partner coverage and lowered placement in every one of the 128 paired comparisons. Net placement differences ranged from $-557$ to $-393$ pupils. Historical-offer visibility stays fixed when we change only the allocation order, so it supplies no additional test in this exercise.
